% Draft for ACM TOSEM (acmart, acmsmall).
% Every number below is copied from results/study2/ (per_source.csv, pairwise.csv,
% rank_by_source.csv, lines_to_reach.csv, complementarity.csv, overlap.csv, ball.csv,
% distance_profile.csv, leak_sensitivity.csv, findings.json). Do not edit numbers by hand;
% re-run src/depgraphs/{study2,analysis2,ball,tokcost,ceiling}.py and re-copy.
\documentclass[acmsmall,review]{acmart}

\usepackage{booktabs}
\usepackage{pgfplots}
\pgfplotsset{compat=1.17}
\newcommand{\TODO}[1]{\textcolor{red}{[TODO: #1]}}
\newcommand{\m}[1]{\texttt{#1}}
\makeatletter
% rows of a table, generated by src/depgraphs/maketables.py
\def\rows#1{\@@input #1.tex }
\makeatother

\acmJournal{TOSEM}
\setcopyright{none}
\acmYear{2026}

\begin{document}

\title{What Belongs in an Agent's Context Window}
\subtitle{Import Structure, Issue Text, and Two Controls That Change the Answer}

\author{Neil Lunavat}
\affiliation{%
  \institution{Pimpri Chinchwad College of Engineering}
  \country{India}}
\email{neillunavat3192@gmail.com}

\begin{abstract}
An agent asked to change one file in an unfamiliar repository cannot read the repository. In
our corpus the median repository is 286{,}385 tokens and 60.5\% of tasks sit in one larger
than a 200{,}000-token window; the files that actually matter are a median 2.98\% of it.
Something must choose. Two signals are available for the choice: the repository's import
structure, and the words of the issue report. Graph-based context tools order by the first.

We compare both, and their combination, on 2{,}412 change tasks from 1{,}162 real pull
requests in 112 Python repositories, scoring each ordering by how much of an evidence-based
answer key it covers per line read. The comparison only has an answer once one says what
``needed'' means, so we grade against two keys and never merge them: the files
\emph{co-edited} in the same pull request, and the files that \emph{define symbols the
change uses}. Ranks below are over all twenty-three orderings we score, the oracle included.

One result survives every control we apply. Reciprocal-rank fusion of a cost-scaled
personalised walk with issue-text-to-path matching is the best non-oracle ordering under
\emph{both} keys on the 702 pull requests that carry both (AUC 0.487 and 0.470). It beats
the best pure lexical ordering by 0.055 and 0.067, and the best pure structural one by 0.097
and 0.039, all Holm-corrected with the repository as the unit.

The two single-signal effects that a less careful reading of our own tables would report do
not survive. Graded on the full corpus, issue-to-path matching beats every purely structural
ordering on the co-edited key (0.434 against 0.389 for the best walk) --- but 40.2\% of
tasks have an issue that names a key file \emph{as code}, and on the rest path matching
falls from fifth to eleventh of the twenty orderings we test against one another, while
\m{hops\_lines}, which reads no issue text at all, takes its place in fifth. Graded on the full corpus, a walk
following imports outward is second on the symbol key (0.491) against path matching's
fifteenth (0.382) --- but the symbol key also covers single-file pull requests, which have no
co-edited key at all, and on a population matched to carry both keys the walk falls to eighth
and is no longer distinguishable from path matching ($p_{\text{Holm}} = 1$). We report both
controls at length because each dissolves an effect that we had ourselves previously written
down as a finding.

What edge direction buys is also narrower than it first appears. Discarding direction costs
nothing --- an undirected walk and an outward one are indistinguishable on the symbol key
($p_{\text{Holm}} = 1$) --- while reversing it is severe ($+0.215$ outward over inward,
$p_{\text{Holm}} < 10^{-17}$). Both figures are over the full corpus, and both hold on the matched
population. The finding is that walking \emph{against} imports is bad,
not that direction is the decisive structural choice. The symbol key is in addition partly
circular with the methods it favours, being built by resolving names with \texttt{jedi},
which follows import edges; we state that at the outset rather than claim a symmetric
discovery.

Two further results bound the problem. The two-hop neighbourhood of the seed, the usual
justification for graph retrieval, holds 87.5\% of the needed files but is also 56\% of the
repository and a median of 21{,}646 lines: it supplies recall, not precision, at a lift of
only 1.57$\times$ over taking the same share of files at random. And the distance to a perfect ordering is
large and measurable without running any model: an oracle puts the complete answer inside an
8{,}000-line budget for 97.5\% of pull requests, whereas the best method we test reaches
58.5\%. Since no agent can use a file it was never given, that gap is an upper bound on
what better selection could add to any system under that budget --- not an estimate of it,
since an oracle reads only key files by construction. We argue that retrieval for code change must be reported per evidence source
\emph{and} per leakage stratum, that filename mentions in issue text are a first-class
confound rather than a footnote, and that a cheap lexical signal belongs in any
graph-based context tool that does not already have one. We make no claim about downstream repair rates;
Section~\ref{sec:protocol} pre-registers the experiment that would test one.
\end{abstract}

\keywords{code search, context retrieval, dependency graphs, software maintenance, empirical
software engineering}

\maketitle

\section{Introduction}

An agent is handed a bug report and a file to change. Before it can write anything it must
decide what else to read, and it cannot simply read everything. In the corpus we study the
median repository is 286{,}385 tokens; 74.8\% of tasks sit in a repository larger than a
128{,}000-token window and 60.5\% in one larger than 200{,}000. Meanwhile the files that
carry the answer are a median 2.98\% of the repository. The reduction available is roughly
34$\times$, and every line spent on an irrelevant file is a line not spent on a relevant
one. Something has to choose, and the choice is the subject of this paper.

A human faces a softer version of the same problem --- attention rather than tokens --- and
we return to that reading in Section~\ref{sec:discussion}. The budget is what makes the
question sharp, so we pose it in those terms throughout.

The dominant answer in tool building is the dependency graph. Aider builds a repository map
from the call and definition graph; RepoGraph, LocAgent and CoSIL construct code graphs and
walk them outward from a starting point. The premise is intuitive and rarely stated: what
you need to read is near, in the graph, to what you are changing.

There is a second signal, and it is the one a human reaches for first. Asked to fix the
login page, a developer thinks of the authentication files --- not because \m{login.py}
imports them, which it may well not, but because the words match. That signal sits in the
issue report.

We are not the first to notice it. SWE-bench's own sparse-retrieval baseline is BM25 with
the issue as the query~\cite{jimenez2024swebench}, Agentless~\cite{xia2025agentless} leans on
file and directory names, and the agentic systems all read the issue text as part of their
loop. What is missing is not the idea that words help. It is a measurement of how far they
carry against structure, on a stated answer key, under a stated budget, and with the issue
corpus's habit of naming the answer controlled for. Every system in Table~\ref{tab:priorkeys}
reports end-to-end accuracy against a single key; none reports that comparison, and none
reports how often its issues name the files it is scored on.

This paper compares the two. The comparison sounds simple and is not, because it has no
answer until one says what it means for a file to have been ``needed''. We use two evidence
sources that can be recovered from the repository without human labelling: the files
\emph{co-edited} in the same pull request, and the files that \emph{define symbols} the
changed code references. Both are defensible. Neither is obviously the right one. They give
opposite answers to the question this paper asks.

\subsection*{Research questions}

\begin{description}
\item[RQ1] How much of a repository is the seed's two-hop neighbourhood, and how much recall
does it buy over a same-sized random sample?
\item[RQ2] Do structural and lexical reading orders rank differently under the two evidence
sources, and by how much?
\item[RQ3] How much does conditioning on the seed file add over a query-independent
repository prior, and does the answer depend on the evidence source?
\item[RQ4] Do the two signals retrieve the same files, or different ones?
\item[RQ5] Does fusing them produce an ordering that is good under both evidence sources?
\item[RQ6] What fraction of tasks could be solved at all under a given budget --- that is,
how much of the gap between current methods and a perfect ordering is headroom rather than a
property of the tasks?
\end{description}

\subsection*{Contributions}

\begin{enumerate}
\item A head-to-head comparison of import-graph proximity and issue-text retrieval as
\emph{reading orders} under a line budget, on 2{,}412 tasks from 1{,}162 real pull requests
in 112 Python repositories, with both signals scored against the same two evidence sources.
\item Evidence that the conclusion of such a comparison is decided as much by the evidence
source and by how the corpus is filtered as by the methods. Across twenty contending
orderings the two keys induce different rankings, but we show that two of the largest
apparent effects are properties of the comparison rather than of the signals: one is
filename leakage from the issue text, the other a population mismatch between the two keys.
All claims are Holm-corrected with the \emph{repository} as the unit, since pull requests
cluster within repositories and treating them as independent inflates significance by many
orders of magnitude.
\item A quantification of what graph proximity actually contributes: recall over half the
repository, at a lift of 1.57$\times$, with 4.2\% precision at two hops.
\item A robustness criterion --- worst-case rank across evidence sources --- under which
three of our five fusions beat every pure signal. The dividing line is not whether path
matching is included, since four of the five include it and one of those still fails. It is
whether the structural input is cost-scaled: the three that fuse from \m{hops\_lines} or
\m{ppr\_und\_pl} succeed, and the two built on the unscaled \m{ppr\_und} do not.
\item A cost accounting in tokens rather than files: a corpus-wide calibration of Python
source (8.50 tokens per line in aggregate, median 7.88), and a method-independent ceiling
--- the share of pull requests whose complete answer fits a given budget at all. An oracle
ordering reaches 97.5\% of pull requests within 8{,}000 lines where the best method we test
reaches 58.5\%, so the shortfall is headroom and not task difficulty.
\item A precise statement of what edge direction buys. Reversing a walk is severely harmful
--- following imports outward beats the same walk run inward by 0.215 AUC on the symbol key,
or 0.123 on the population matched to both keys ($p_{\text{Holm}} < 10^{-14}$ in both cases,
repository as the unit) --- but \emph{discarding} direction costs nothing, because the
undirected walk common in this literature is indistinguishable from the outward one
($p_{\text{Holm}} = 1$ on the full corpus, $0.42$ on the matched one). On the co-edited key direction is not merely
unhelpful but harmful: all four directed variants rank seventeenth to twentieth, below a
same-directory baseline that uses no graph at all, while all four undirected walks rank tenth
to fourteenth, above it.

\item Two controls that each dissolve an effect we had previously reported as a finding, and
which we recommend as standard practice. A \emph{leakage} control that classifies whether an
issue names a key file as code rather than merely sharing a word with it: 40.2\% of tasks do,
and on the rest issue-to-path matching falls from fifth to eleventh of the twenty contenders
on the co-edited key. A
\emph{population} control that restricts both keys to the 702 pull requests carrying both,
since a co-edited key requires a multi-file pull request while a symbol key does not: on the
matched population the best directed walk falls from second to eighth on the symbol key and
is no longer distinguishable from path matching.
\item A replication package with the graphs, tasks, issue-token bags, token counts,
orderings, generated tables and analysis scripts.
\end{enumerate}

\section{Background and Related Work}

\paragraph{Dependency structure as a guide to change.} Using program structure to decide what
a change touches is old. Impact analysis~\cite{bohner1996impact} propagates from a changed
element along dependencies; Robillard~\cite{robillard2008topology} ranks program elements by
the topology of their dependencies; feature-location surveys~\cite{dit2013featurelocation} catalogue
structural, textual and dynamic techniques and note that they are usually evaluated in
isolation. Ko et al.~\cite{ko2006exploratory} observed developers navigating largely by
identifier names and search rather than by following dependencies.

\paragraph{Change history.} Zimmermann et al.~\cite{zimmermann2004mining} and Ying et
al.~\cite{ying2004predicting} mine co-change from version histories to suggest files that
change together. Our co-edited evidence source is the same notion of relatedness, used as
ground truth rather than as a method.

\paragraph{Graph-based context for agents.} SWE-bench~\cite{jimenez2024swebench} and its
successors~\cite{badertdinov2025swerebench,zhang2025swebenchlive} made repository-scale
change tasks a standard benchmark, and a family of systems now assembles context by walking
a code graph: Aider's repository map~\cite{aider2023repomap},
RepoGraph~\cite{ouyang2025repograph}, LocAgent~\cite{chen2025locagent} and
CoSIL~\cite{jiang2025cosil}. Agentless~\cite{xia2025agentless} is a notable counterexample that
leans on file and directory names.

What these systems share is the ground truth they are judged against.
Table~\ref{tab:priorkeys} records it for every system whose evaluation section we read in
the published PDF; each entry is quoted or paraphrased from that section, not from an
abstract or a secondary description.
Where localization accuracy is reported at all, the target is the set of files or functions
touched by the gold patch --- the SWE-bench convention, and exactly our \m{co\_edited}
source. That is the evidence source on which, in our results, structural proximity is the
weaker of the two signals.

\begin{table*}
\caption{What repository-context systems are graded against, read from each paper's own
evaluation section (2026-09-23). The upper block is the patch-based convention; the lower
block is the concurrent work that departs from it.}
\label{tab:priorkeys}
\footnotesize
\setlength{\tabcolsep}{4pt}
\begin{tabular}{llp{8.8cm}}
\toprule
System & Venue & Ground truth used for evaluation \\
\midrule
\multicolumn{3}{l}{\emph{Files touched by the gold patch}} \\
LocAgent~\cite{chen2025locagent} & ACL 2025 & ``The affected files or functions \ldots as
identified in the patches, are considered the target locations'' \\
CoSIL~\cite{jiang2025cosil} & ASE 2025 & File- and function-level accuracy on SWE-bench Lite
and Verified, i.e.\ the patched locations \\
Agentless~\cite{xia2025agentless} & FSE 2025 & Correct if the edit covers a superset of the
ground-truth patch locations, at line, function and file level \\
LingmaAgent~\cite{ma2025lingmaagent} & FSE 2025 & Fault locations extracted from the
developer patch, at file and function level \\
SWE-Fixer~\cite{xie2025swefixer} & ACL 2025 & Precision and recall against the defective
files of the resolving pull request; BM25 top-30 then a learned reranker \\
KGCompass~\cite{yang2025kgcompass} & 2025 & Files and code elements modified by the
ground-truth patch, scored by Jaccard overlap \\
SWE-Adept~\cite{he2026sweadept} & 2026 & Correct only if \emph{all} locations modified by
the ground-truth patch appear in the top-$k$ \\
SWERank~\cite{reddy2025swerank} & 2025 & Every function modified in the pull request is a
positive \\
CodeScout~\cite{sutawika2026codescout} & 2026 & Reuses LocAgent's patch-processing scripts:
modified files, modules and functions \\
RepoGraph~\cite{ouyang2025repograph} & ICLR 2025 & ``The percentage of problems where
the edit locations match the ground truth patch'', at file, function and line level \\
\midrule
\multicolumn{3}{l}{\emph{Something other than the gold patch alone}} \\
GraphLocator~\cite{liu2025graphlocator} & 2025 & Locations extracted from the fix patch,
then dual-annotated and cross-reviewed because patch-derived ground truth ``contains
missing locations'' \\
SWE-Explore~\cite{zhang2026sweexplore} & 2026 & File and line intervals that successful
agents actually \emph{read}, intersected across several verified trajectories \\
ContextBench~\cite{li2026contextbench} & 2026 & Human-annotated context traced outward from
the patch regions, pruned by a second annotator, validated by an LLM patching from it alone \\
AgentRetrievalBench~\cite{qin2026agentretrievalbench} & 2026 & A file is relevant if reading
it helps the agent's next step; one of four task types is ripple files beyond an anchor \\
CodeR~\cite{chen2024coder} & 2024 & End-to-end on SWE-bench, but its separate localization
study uses golden locations marked by the authors \\
\bottomrule
\end{tabular}
\end{table*}

\paragraph{Concurrent work is questioning the key.} The patch convention is under active
revision, by at least four efforts that appeared while this study was running, and we want
to be exact about how close they come. SWE-Explore~\cite{zhang2026sweexplore} derives
line-level ground truth ``from independent agent trajectories that successfully solved the
same issue, distilling the specific code regions their solution paths actually consulted'',
and scores a ranked list ``under a fixed line budget'' --- both the key and the cost model
are close relatives of ours. ContextBench~\cite{li2026contextbench} annotates gold context
by hand, 522{,}115 lines over 1{,}136 tasks, and measures recall, precision and efficiency
rather than end-to-end resolution. AgentRetrievalBench~\cite{qin2026agentretrievalbench}
defines a file as relevant ``when reading it would materially help the agent take the next
correct step'', evaluated under a top-$k$ or token budget; one of its four task types,
ripple files beyond an anchored edit, is our co-edited source by another name. And
GraphLocator~\cite{liu2025graphlocator} keeps the patch key but dual-annotates it, on the
stated grounds that patch-derived ground truth misses locations.

We take this as corroboration of the premise rather than as a threat to it: several groups
independently concluded in the same year that the patch key is inadequate. What none of
them reports is the consequence we measure. Each proposes a better key and evaluates
methods against it; we hold two keys already recoverable from existing benchmark data, and
show that the choice between them \emph{reverses the ranking of the methods themselves},
with every individual reversal significant under Holm correction. The contribution is not the observation that keys differ. It is the
measurement of what the field's default key costs the signal the field has invested in.

\paragraph{Retrieval.} BM25~\cite{robertson2009bm25} remains the reference lexical ranking
function; we use it unchanged, varying only the query.

\section{Study Design}

\label{sec:design-freeze}
The corpus, the graph builder, the task definition and the coverage--cost metric were fixed
and tagged in version control before any method was run; the present study adds method
families and analyses on top of that frozen base. Changes made after that tag, and claims we
have retracted, are recorded in the replication package rather than silently absorbed.

The freeze covers the corpus, the graphs, the task definition and the metric. It does not
cover the method families: the fusion orderings, the directed walks and the lexical methods
queried by issue text were all added afterwards, in response to what the first results
showed. Fusion is the paper's main recommendation and it is an \emph{exploratory} result, not
a pre-registered one. We say so here rather than let the word ``pre-registered'' cover more
than it earns.

\subsection{Repositories and tasks}

We draw pull requests from SWE-bench, SWE-bench-Live, SWE-Gym, SWE-bench Pro, Loc-Bench and
SWE-rebench, keeping Python repositories with a permissive licence. A pull request that
edits $k$ existing non-test source files that are nodes of the import graph yields $k$
tasks: each edited file in turn becomes the \emph{seed}, and the task is to order the
remaining files of the repository. This gives 3{,}123 tasks from 1{,}873 pull requests in
113 repositories. Of these, 2{,}412 tasks --- 1{,}162 pull requests in 112 repositories ---
have at least one key file inside the graph once the seed is removed, and they are the ones
scored here; the remaining 711 tasks have an empty key under both sources and are reported
but not scored. Those 711 tasks come from 711 distinct pull requests, and 706 of them change
a single file: with no second edited file there is no co-edited key, and the symbol key is
empty whenever the changed code resolves to nothing outside itself --- standard library,
third-party packages, or names \texttt{jedi} could not resolve. That is 61.2\% of all
single-file tasks, which is a real limitation of the symbol key and not a filtering choice
of ours. Within the scored set, 1{,}958 tasks (708 pull requests) have a co-edited
key and 2{,}341 (1{,}156 pull requests) a symbol key.

The median repository contributes 131 graph nodes and the median key is 3 files totalling
1{,}515 lines, so the task is to find a handful of files among roughly a hundred.

\subsection{Import graphs}

Graphs are built per base commit with a custom \m{ast} walker over the repository at that
commit. Nodes are source files; an edge $a \rightarrow b$ means $a$ imports a module
resolving to $b$. Package \m{\_\_init\_\_} files are not inserted as intermediaries,
\m{from X import Y} resolves to the module, edges are unweighted, imports inside
\m{TYPE\_CHECKING} blocks and inside function bodies are kept, and test files are excluded.
Per-file line counts come from the same pass and are the cost unit throughout.

\subsection{Two evidence sources}

\begin{description}
\item[\m{co\_edited}] The other non-test source files changed by the same pull request. This
is what a developer in fact touched, and it is independent of the import graph.
\item[\m{symbol}] Files that define names the changed code references, resolved with a static
analyser. This is what the change depends on.
\end{description}

\noindent Symbol resolution uses \texttt{jedi}'s \texttt{goto} with
\texttt{follow\_imports=True}, which reaches a name's defining file by traversing import
statements. This is worth stating plainly because it bounds what the symbol key can show:
the key is constructed by walking the import graph, so a method that walks the import graph
is being scored against a target built from the same substrate. We report it because it is
the natural static notion of ``what this code depends on'', not because it is a neutral
arbiter between the two method families.

\noindent The seed file is removed from both. We deliberately do not merge them: the whole
question is whether the choice matters. It does, and the union is not a neutral compromise
--- the symbol key contributes 7{,}675 file references against 4{,}458 for co-editing, so a
union largely inherits the structural source's verdict.

A third source --- files actually executed by the tests that the fix repairs --- is defined
in our protocol and was collected for 113 pull requests before we stopped. We exclude it
rather than report it: a source covering under a tenth of the corpus cannot be compared
against two that cover all of it, and reporting it alongside them would invite exactly the
conflation this paper argues against. The partial collection is in the replication package
so that it is visible rather than quietly dropped, and completing it is the obvious
extension.

\subsection{Reading orders}

Every method receives the graph at the base commit, the seed file, per-file identifier bags
and the issue text, and returns an ordering of all other nodes. Table~\ref{tab:methods}
lists them by family. Identifier bags are built by splitting every identifier in the file on
underscores and camel-case boundaries and lower-casing; the issue text is tokenised
identically, so text is matched against code over one vocabulary. BM25 is used with the
standard $k_1 = 1.5$, $b = 0.75$ and differs between methods only in the query. Fusion is
reciprocal-rank fusion with $k = 60$.

Two points of fairness are worth stating. The structural methods are \emph{given} a correct
answer for free --- the seed is one of the files the pull request changed. The issue-text
methods are given no seed at all. And the metric charges by lines read, so we include
cost-scaled variants (suffix \m{\_pl}, score divided by file length) to separate ``ranks the
right files highly'' from ``happens to prefer short files''.

\begin{table}
\caption{Reading orders, by family. Structural methods use the seed and the graph; global
methods use neither the seed nor the issue; lexical methods use no graph.}
\label{tab:methods}
\small
\begin{tabular}{llp{7.2cm}}
\toprule
Family & Method & Ordering \\
\midrule
structure & \m{bfs\_und} & breadth-first from the seed, undirected \\
 & \m{hops\_lines} & by (hop distance from seed, file length) \\
 & \m{ppr\_und} & personalised PageRank from the seed, undirected \\
 & \m{ppr\_und\_pl} & the same, divided by file length \\
 & \m{ppr\_out} & the same walk along import edges: what the seed depends on \\
 & \m{ppr\_out\_pl} & the same, divided by file length \\
 & \m{ppr\_in} & the same walk against import edges: what depends on the seed \\
 & \m{ppr\_in\_pl} & the same, divided by file length \\
\midrule
global & \m{pagerank} & PageRank of the whole graph; ignores seed and issue \\
 & \m{indegree} & number of files importing this one \\
\midrule
lexical & \m{bm25\_seed} & BM25, query = identifiers of the seed file \\
 & \m{path\_seed} & BM25 over paths, query = the seed's path tokens \\
 & \m{bm25\_issue} & BM25, query = identifiers of the issue text \\
 & \m{bm25\_issue\_pl} & the same, divided by file length \\
 & \m{path\_issue} & BM25 over file paths, query = issue tokens \\
\midrule
fusion & \m{rrf\_hops\_path} & RRF of \m{hops\_lines} and \m{path\_issue} \\
 & \m{rrf\_hops\_issue\_path} & RRF of \m{hops\_lines}, \m{bm25\_issue}, \m{path\_issue} \\
 & \m{rrf\_pprpl\_issue\_path} & RRF of \m{ppr\_und\_pl}, \m{bm25\_issue}, \m{path\_issue} \\
 & \m{rrf\_ppr\_issue} & RRF of \m{ppr\_und} and \m{bm25\_issue} \\
 & \m{rrf\_ppr\_issue\_path} & RRF of \m{ppr\_und}, \m{bm25\_issue}, \m{path\_issue} \\
\midrule
reference & \m{same\_dir} & files in the seed's directory first \\
 & \m{random} & uniformly shuffled: the floor \\
 & \m{oracle} & key files only, cheapest first: the ceiling \\
\bottomrule
\end{tabular}
\end{table}

\subsection{Metric}

An ordering is read from the top until a line budget is exhausted; coverage is the fraction
of the key that has been read. We evaluate at budgets
$\{250, 500, 1000, 2000, 4000, 8000, 16000, 32000\}$ and report the AUC, the mean coverage
over those eight doublings, so that a method must be good at both tight and generous
budgets. We also report the more directly interpretable quantity: the number of lines that
must be read before half the key has been covered.

Orderings with ties are re-drawn three times per task with distinct seeds and the coverage
curves averaged.

The budget is denominated in lines because the pre-registered design fixed it that way at
\texttt{freeze-1}, and we do not change a frozen metric after seeing results. Lines are
nevertheless the wrong unit for a context window, so we measure the conversion rather than
assume it: across the 121{,}075 distinct file versions in the corpus, tokenised with the
\texttt{cl100k\_base} byte-pair vocabulary, Python source runs to 8.50 tokens per line in
aggregate and a median of 7.88 per file, with a 5th--95th percentile range of 5.70--10.90.
Absolute counts shift by a small constant factor between tokenizer families; the ratios we
report do not depend on the choice. A budget of 8{,}000 lines is therefore roughly 68{,}000
tokens, and the eight budgets span roughly 2{,}100 to 272{,}000 tokens --- from a
fraction of a small context window to more than a large one.

\subsection{Statistics}

Task AUCs are averaged within a pull request first, so a change touching many files does not
count many times. Paired comparisons use the Wilcoxon signed-rank test with the
Vargha--Delaney $A_{12}$ effect size, Holm-corrected within each family of comparisons.
Confidence intervals are 2{,}000-resample bootstraps resampling \emph{repositories}, not
tasks, so that a few large repositories cannot manufacture significance.

The paired tests are reported with the \textbf{repository} as the unit: the pull-request
scores are averaged within a repository and the test runs over the 112 repositories. Pull
requests are not independent draws --- they cluster in repositories that differ in size,
layout convention and issue-writing culture --- and a test that ignores the clustering
counts the same evidence many times. The difference is not cosmetic. The same comparison
that yields $p < 10^{-100}$ over 1{,}156 pull requests yields $p < 10^{-20}$ over 112
repositories. The effect sizes are essentially unchanged, which is the point: the conclusions
survive, and only the arithmetic of the $p$-values was ever overstated. We report the
pull-request-level tests alongside, in the replication package, so the two can be compared.

Because a co-edited key requires a pull request that touches more than one graph file while
a symbol key does not, the two keys do not cover the same pull requests. Wherever we compare
across keys we therefore report the comparison twice: over each key's full population, and
over the 702 pull requests that carry both (Section~\ref{sec:shared}).

\section{Results}

\subsection{RQ1: the two-hop neighbourhood is half the repository}

``The files you need are within two hops'' is true here: 86.7\% of symbol-key files and
88.1\% of co-edited files lie within two undirected hops of the seed
(Table~\ref{tab:distance}). Taken alone this reads as strong evidence for graph retrieval.

The other half of the statistic undoes it. In the median task the two-hop ball holds 68 of
the repository's 131 files and 21{,}646 lines --- 56\% of the files on average, and more
than five times a 4{,}000-line budget. Its precision is 4.2\%. Against a random sample of
the same size it lifts recall by a factor of 1.57.

The two-hop neighbourhood is therefore not a filter. It is a statement that the files you
need are somewhere in about half the repository --- which is exactly the regime in which the
\emph{ordering} decides everything and the graph has already said all it has to say.

\begin{table}
\caption{Where the key files sit relative to the seed, and how big the two-hop ball is.
Hop shares are pooled proportions over all key-file references; the ball columns are medians
over 2{,}412 tasks.}
\label{tab:distance}
\small
\begin{tabular}{lrrrr}
\toprule
Key files at hop distance & 1 & 2 & 3+ & unreachable \\
\midrule
\m{co\_edited} & 39.0\% & 49.1\% & 9.1\% & 2.8\% \\
\m{symbol} & 45.3\% & 41.4\% & 9.1\% & 4.2\% \\
\midrule
\multicolumn{5}{l}{\emph{The two-hop ball:} 68 files, 21{,}646 lines, 56\% of the
repository} \\
\multicolumn{5}{l}{\quad recall 87.5\%, precision 4.2\%, recall lift over a same-sized
random sample 1.57$\times$} \\
\multicolumn{5}{l}{\emph{The key:} 3 files, 1{,}515 lines} \\
\bottomrule
\end{tabular}
\end{table}

\subsection{RQ2: the evidence source decides the winner}

Table~\ref{tab:main} is the central result. Under the co-edited key the top five orderings
are four fusions and one lexical method; the best purely structural ordering is tenth, and
the query-independent priors are twenty-first and twenty-second of twenty-three, below
\m{same\_dir} and barely above random. Under the symbol key the table inverts: the top four
are three structural orderings and one fusion, and the best lexical ordering is fifteenth.

The individual reversals are large and significant after Holm correction with the repository
as the unit. Issue-text path matching beats cost-scaled personalised PageRank by 0.059 AUC
on co-edited files ($A_{12} = 0.403$ for the walk, $p_{\text{Holm}} < 10^{-4}$) and loses to
it by 0.100 on symbol files ($A_{12} = 0.671$, $p_{\text{Holm}} < 10^{-8}$).

We rest the reversal on those paired tests rather than on a rank correlation. The
correlation is low --- Spearman $\rho = 0.24$ over the twenty contenders ($p = 0.31$), and
$\rho = 0.37$ over all twenty-two non-oracle orderings ($p = 0.09$) --- but with twenty
points a non-significant $\rho$ is weak evidence of anything, and its value moves with which
methods are included. The Holm-corrected pairwise reversals do not have that weakness.

Two things about Table~\ref{tab:main} are \emph{not} safe to read off it, and the next two
subsections are about them. The two columns are computed over different pull requests, and
the left column is computed over a corpus in which many issues name the answer outright.
Each of those turns out to account for one of the two halves of the reversal.

Figure~\ref{fig:curves} shows the same reversal as coverage curves. Table~\ref{tab:lines}
states it in the unit a context budget is actually spent in: to reach half the co-edited
files a reader following cost-scaled personalised PageRank consumes a median of 6{,}155
lines, against 4{,}838 following issue-text path matching and 3{,}342 following the best
fusion. On symbol files the same three numbers are 3{,}483, 6{,}638 and 3{,}957 --- the walk
and the words trade places, and the fusion is close to the better of the two in both.

\subsection{The reversal is smaller on a matched population}
\label{sec:shared}

A co-edited key exists only for a pull request that changes more than one file in the graph.
A symbol key does not need that, and 444 of the 1{,}156 pull requests with a symbol key
change a single file. The two columns of Table~\ref{tab:main} are therefore scored on
different samples, and single-file changes are exactly the case where a seed's import
neighbourhood is most of what there is to find. The comparison has to be repeated on a
population that carries both keys.

Table~\ref{tab:shared} does that, on the 702 pull requests that have both. On the co-edited
key nothing moves: no method changes rank by more than one place. On the symbol key the
structural family loses most of its advantage. Cost-scaled personalised PageRank along
import edges falls from second to eighth (0.491 to 0.417), the undirected form from third to
fifth, and \m{hops\_lines} from fourth to sixth, while issue-to-path matching rises from
fifteenth to ninth. The three top places under both keys are then held by the same three
fusions.

The statistical claim moves with the ranks. On the full corpus a walk along import edges
beats path matching on the symbol key by 0.101 ($p_{\text{Holm}} < 10^{-8}$, repository as
the unit). On the matched population the same comparison gives $+0.008$ and
$p_{\text{Holm}} = 1$; for \m{hops\_lines} against path matching it gives $+0.019$ and
$p_{\text{Holm}} = 1$. We therefore do not claim that structural proximity beats issue-text
retrieval on the symbol key. On a like-for-like sample the two are indistinguishable, and
what remains of the reversal is the co-edited half --- which the next control also qualifies.

\begin{table}
\caption{The two keys re-scored on the 702 pull requests that carry both, top ten under
each. $\Delta$ is the change in rank against the full population of
Table~\ref{tab:main}; a positive number means the method does worse once the samples are
matched.}
\label{tab:shared}
\small
\begin{tabular}{llcrc}
\toprule
Method & Fam. & AUC & \# & $\Delta$ \\
\midrule
\rows{tab_shared}
\bottomrule
\end{tabular}
\end{table}

\begin{table}
\caption{Mean AUC by evidence source, pull request as the unit, with bootstrap 95\%
confidence intervals over repositories. The same twenty-three orderings, ranked twice.}
\label{tab:main}
\footnotesize
\setlength{\tabcolsep}{3.5pt}
\begin{tabular}{rllc@{\hskip 1.2em}rllc}
\toprule
\multicolumn{4}{c}{\m{co\_edited} (708 PRs)} & \multicolumn{4}{c}{\m{symbol} (1{,}156 PRs)} \\
\cmidrule(r){1-4}\cmidrule(l){5-8}
\# & Method & Fam. & AUC [CI] & \# & Method & Fam. & AUC [CI] \\
\midrule
\rows{tab_main}
\bottomrule
\end{tabular}
\end{table}

\begin{figure}
\centering
\begin{tikzpicture}
\begin{axis}[width=0.46\textwidth, height=5.4cm, xmode=log, log basis x=2,
  xlabel={lines read}, ylabel={coverage of the key}, title={\m{co\_edited}},
  legend pos=north west, legend style={font=\tiny, draw=none, fill opacity=0.7,
  text opacity=1}, ymin=0, ymax=1, grid=major,
  tick label style={font=\small}, label style={font=\small}, title style={font=\small}]
\addplot table[x=budget,y=oracle] {fig_co_edited.dat}; \addlegendentry{oracle}
\addplot table[x=budget,y=rrf_hops_path] {fig_co_edited.dat}; \addlegendentry{rrf\_hops\_path}
\addplot table[x=budget,y=path_issue] {fig_co_edited.dat}; \addlegendentry{path\_issue}
\addplot table[x=budget,y=ppr_und_pl] {fig_co_edited.dat}; \addlegendentry{ppr\_und\_pl}
\addplot table[x=budget,y=pagerank] {fig_co_edited.dat}; \addlegendentry{pagerank}
\addplot[black,dashed] table[x=budget,y=random] {fig_co_edited.dat}; \addlegendentry{random}
\end{axis}
\end{tikzpicture}\hfill
\begin{tikzpicture}
\begin{axis}[width=0.46\textwidth, height=5.4cm, xmode=log, log basis x=2,
  xlabel={lines read}, ylabel={}, title={\m{symbol}},
  legend pos=north west, legend style={font=\tiny, draw=none, fill opacity=0.7,
  text opacity=1}, ymin=0, ymax=1, grid=major,
  tick label style={font=\small}, label style={font=\small}, title style={font=\small}]
\addplot table[x=budget,y=oracle] {fig_symbol.dat}; \addlegendentry{oracle}
\addplot table[x=budget,y=rrf_hops_path] {fig_symbol.dat}; \addlegendentry{rrf\_hops\_path}
\addplot table[x=budget,y=path_issue] {fig_symbol.dat}; \addlegendentry{path\_issue}
\addplot table[x=budget,y=ppr_und_pl] {fig_symbol.dat}; \addlegendentry{ppr\_und\_pl}
\addplot table[x=budget,y=pagerank] {fig_symbol.dat}; \addlegendentry{pagerank}
\addplot[black,dashed] table[x=budget,y=random] {fig_symbol.dat}; \addlegendentry{random}
\end{axis}
\end{tikzpicture}
\caption{Mean coverage against lines read. The same five orderings change places between the
two evidence sources: issue-text path matching tracks the best fusion on co-edited files and
falls below the query-independent prior on symbol files, while cost-scaled personalised
PageRank does the reverse.}
\label{fig:curves}
\end{figure}

\begin{table}
\caption{Median lines that must be read before half, then all, of the key is covered. Lower is better.}
\label{tab:lines}
\small
\begin{tabular}{lrr@{\hskip 1.5em}rr}
\toprule
& \multicolumn{2}{c}{to half the key} & \multicolumn{2}{c}{to all of it} \\
\cmidrule(r){2-3}\cmidrule(l){4-5}
Method & \m{co\_edited} & \m{symbol} & \m{co\_edited} & \m{symbol} \\
\midrule
\rows{tab_lines}
\bottomrule
\end{tabular}
\end{table}

\subsection{Direction: reversing a walk is costly, discarding it is free}
\label{sec:direction}

The walks above are undirected, as they are in most of this literature. An import edge has a
direction, so we add both orientations of the personalised walk, cost-scaled and not.

The result is asymmetric, and the asymmetry is the finding. Running the walk \emph{along}
import edges --- towards what the seed depends on --- places \m{ppr\_out\_pl} second of
twenty-three on the symbol key at 0.491. Running it \emph{against} them collapses to
\m{ppr\_in\_pl} at 0.277, twenty-first. The like-for-like gap is 0.123 on the matched
population and 0.215 on the full one, in both cases the largest separation between two
structural orderings anywhere in our results ($p_{\text{Holm}} < 10^{-14}$, repository as
the unit).

But the undirected walk loses almost nothing. On the symbol key \m{ppr\_und\_pl} scores
0.489 against \m{ppr\_out\_pl}'s 0.491, a difference of 0.002 that is not significant on the
full corpus ($p_{\text{Holm}} = 1$) and reverses in sign, slightly, on the matched one. So
the practice this literature has settled on --- throw the direction away --- is not the
mistake it might appear to be from the spread of the table. Walking backwards is the mistake.

On the co-edited key direction is worse than useless. All four directed variants rank
seventeenth to twentieth, below \m{same\_dir}, a baseline that orders files by directory and
consults no graph at all; all four undirected walks rank tenth to fourteenth, above it.
Committing to an orientation costs 0.044 to 0.061 AUC there ($p_{\text{Holm}} < 10^{-3}$).

The explanation is the same one that qualifies the symbol key throughout. That key is built
by resolving names with \texttt{jedi}, which follows imports in exactly one direction, so a
walk in that direction is measuring something close to the key's own construction. The
co-edited key records a human decision about what belonged in one change, and no orientation
of the import graph recovers it well.

\subsection{RQ3: what the seed is worth depends on the source too}

A query-independent PageRank prior --- a ranking that never looks at the seed or the issue,
and simply says which files matter in this repository --- reaches 0.410 on symbol files
against 0.489 for the best cost-scaled personalised walk, which does look at the seed. That
is 84\% of the personalised score, and far above random (0.248; $A_{12} = 0.758$,
$p_{\text{Holm}} < 10^{-16}$). On that evidence source, much of what the personalised walk
delivers is a repository prior.

Under the co-edited key the same prior scores 0.271 against a random floor of 0.248. The
difference is significant but negligible ($A_{12} = 0.552$, $p_{\text{Holm}} = 0.001$), and
the personalised walk beats it by 0.101 ($A_{12} = 0.666$). Knowing which file is being changed therefore
matters a great deal for predicting what else was changed, and rather little for predicting
what the code calls, where knowing the repository's hubs is nearly as good.

This is worth stating plainly, because the opposite claim is easy to make from a single
source: whether ``the graph knows what you are working on'' appears true or false is itself
a function of the answer key.

\subsection{RQ4: the two signals name different files}
\label{sec:rq4}

The reversal is not two methods disagreeing about the order of one list. At $k=10$ the
top-10 sets of personalised PageRank and issue-text BM25 have a mean Jaccard overlap of
0.187, and of personalised PageRank and issue-text path matching only 0.102; at $k=50$ these
rise to 0.508 and 0.412. For comparison, personalised PageRank and the query-independent
prior overlap at 0.364 at $k=10$: a seed-conditioned walk resembles the repository prior
more than it resembles a lexical ranking.

Table~\ref{tab:comp} follows individual key files. At $k=20$, 24.4\% of co-edited files are
found only by the lexical family and 9.6\% only by the structural family; on symbol files the
asymmetry flips (19.9\% structure-only against 16.5\% words-only). Under either source about
a sixth of key files are found by neither family at $k=20$, which is where the remaining
headroom to the oracle lives.

\begin{table}
\caption{How each key file is reached within the top $k$ of the structural family
(\m{ppr\_und}, \m{hops\_lines}) and the lexical family (\m{bm25\_issue}, \m{path\_issue}).}
\label{tab:comp}
\small
\begin{tabular}{llrrrr}
\toprule
Source & $k$ & both & structure only & words only & neither \\
\midrule
\m{co\_edited} & 10 & 32.2\% & 11.4\% & 30.0\% & 26.4\% \\
\m{co\_edited} & 20 & 50.7\% & 9.6\% & 24.4\% & 15.3\% \\
\m{co\_edited} & 50 & 76.7\% & 5.9\% & 13.4\% & 4.0\% \\
\m{symbol} & 10 & 28.5\% & 23.3\% & 20.4\% & 27.8\% \\
\m{symbol} & 20 & 47.4\% & 19.9\% & 16.5\% & 16.2\% \\
\m{symbol} & 50 & 72.9\% & 11.5\% & 9.8\% & 5.9\% \\
\bottomrule
\end{tabular}
\end{table}

\subsection{RQ5: fusion is the only ordering that survives both sources}

If the evidence source is not known in advance --- and for a tool shipped to users it is not
--- the quantity to optimise is the worst rank a method attains across sources.
Table~\ref{tab:worst} ranks by it. Three of the five fusions beat every pure signal. The
two that do not --- \m{rrf\_ppr\_issue} and \m{rrf\_ppr\_issue\_path} --- are not
distinguished by whether they include path matching, since four of the five do and one
of those still fails. They are the two built on the unscaled \m{ppr\_und}. The three that
survive all fuse from a cost-scaled structural input, either \m{hops\_lines} or
\m{ppr\_und\_pl}. Dividing by file length before fusing is what makes the difference. The best,
\m{rrf\_hops\_path}, is nothing more than reciprocal-rank fusion of hop-distance ordering
with issue-tokens-against-file-paths, and is fourth under one source and fifth under the
other.

The insurance is close to free. Against pure structure (\m{hops\_lines}) it gains 0.067 AUC
on co-edited files ($A_{12} = 0.602$, $p_{\text{Holm}} < 10^{-11}$) and gains 0.003 on
symbol files ($p_{\text{Holm}} = 1$); against \m{ppr\_und\_pl} it gains 0.087
($p_{\text{Holm}} < 10^{-13}$) and loses 0.008 ($p_{\text{Holm}} = 1$). The lexical signal adds a large amount where structure is weak and costs
nothing where structure is strong.

\begin{table}
\caption{Worst rank attained across the two evidence sources. Ranks are those of
Table~\ref{tab:main}, over all twenty-three orderings with the oracle included, so the
best attainable worst rank here is 2. Lower is better. The three fusions built on a
cost-scaled structural input beat every pure signal; the two built on the unscaled walk
do not.}
\label{tab:worst}
\small
\begin{tabular}{llccc}
\toprule
Method & Fam. & \m{co\_edited} & \m{symbol} & worst \\
\midrule
\rows{tab_worst}
\bottomrule
\end{tabular}
\end{table}

\subsection{RQ6: what a budget rules out, and what it does not}
\label{sec:ceiling}

A coverage number says how much of the key a method found. It does not say whether finding
all of it was possible. If the needed files together exceed the budget, no ordering of them
fits and no downstream system can succeed, however capable: a model cannot use a file it was
never given. That bound is a property of the task and the budget alone, so it can be
measured once and applied to every method at once --- and, usefully here, without running a
model.

Table~\ref{tab:ceiling} reports it. The complete answer fits a 32{,}000-token budget for
90.3\% of pull requests under the co-edited key and 87.3\% under the symbol key, and a
64{,}000-token budget for over 95\% of both. The task is not, in general, impossible at
realistic budgets.

We are careful about what this licenses. The oracle reads only key files, so that it fits a
small key into a small budget is close to arithmetic: the ceiling shows the keys are small,
not that the remaining distance is reachable from the signals available. What it does rule
out is the competing explanation --- that current methods score as they do because the
answers do not fit --- and that is worth ruling out, because it is the explanation a reader
reaches for first. The gap between 58.5\% and 97.5\% is therefore an upper bound on what
better selection could add, not an estimate of it.

\begin{table}
\caption{Share of pull requests whose \emph{entire} key fits within a token budget. This is
a ceiling: no ordering can exceed it, and no agent can act on files outside it.}
\label{tab:ceiling}
\small
\begin{tabular}{rrr}
\toprule
Token budget & co\_edited (\%) & symbol (\%) \\
\midrule
\rows{tab_ceiling}
\bottomrule
\end{tabular}
\end{table}

Against that ceiling, current orderings fall a long way short. Table~\ref{tab:full} counts
the pull requests each method takes to \emph{complete} coverage at 8{,}000 lines, a budget
at which an oracle succeeds on 97.5\% of them. The best method on the co-edited key,
\m{rrf\_hops\_issue\_path}, reaches 58.5\%; the best on the symbol key, \m{hops\_lines},
reaches 51.0\%. Roughly forty points of the gap is therefore attributable to the ordering
rather than to the budget or the task. The ceiling is partly definitional --- an oracle that
reads only key files must win --- but its value is the converse reading: at 8{,}000 lines
the budget is almost never the binding constraint, so a failure to cover the key is a
failure of the ordering and not of the budget.

\begin{table}
\caption{Share of pull requests taken to \emph{complete} coverage within 8{,}000 lines. The
oracle row is the ceiling of Table~\ref{tab:ceiling} expressed in the frozen unit. The
reversal of Table~\ref{tab:main} reappears under this stricter criterion.}
\label{tab:full}
\small
\begin{tabular}{lrr}
\toprule
Ordering & co\_edited (\%) & symbol (\%) \\
\midrule
\rows{tab_full}
\bottomrule
\end{tabular}
\end{table}

Two things are worth noting. The reversal survives the stricter criterion:
\m{bm25\_issue} loses 14.9 points and \m{path\_issue} 11.5 points moving from the co-edited
key to the symbol key, while the outward walks gain the most (\m{ppr\_out} 15.2 points,
\m{ppr\_out\_pl} 14.7) and \m{pagerank} gains 11.5. The gain is not a property of
the structural family as such: the two inward walks lose 11.2 and 11.7 points, more than
\m{path\_issue} does. It is a property of walking the direction the key was built along.
And the fusion orderings again sit at or near the top of both columns, which is the same
conclusion RQ5 reached by AUC, reached here by a measure that ignores partial credit
entirely.

\subsection{Sensitivity: are the lexical methods just copying paths out of the issue?}
\label{sec:leak}

Issue reports often contain stack traces, and a stack trace names files. If the lexical
advantage came from reports quoting the paths of the files that had to change, it would be
an artefact of the benchmark rather than a property of language.

Largely it does, and the effect is big enough that we report it as a result rather than as a
reassurance.

The full path of a key file appears verbatim in the issue text for 12.6\% of tasks, and a
bare file stem for 48.5\%. Path BM25 tokenises a path into lowercased subtokens, so a stem is
enough to match and the full-path control is far too weak. But a bare stem is also too
strong a criterion on its own: an issue about ``models'' and a key file called
\texttt{models.py} share a word because a project's prose and its filenames draw on one
domain vocabulary, and using that overlap is retrieval, not leakage. We therefore classify
every stem occurrence by how it is written. A mention is \emph{explicit} if it appears as
code --- \texttt{core.py}, \texttt{simbad/core}, a dotted module path, inside backticks, or
inside a fenced or indented block --- and \emph{incidental} if the stem occurs only as an
ordinary English word. Explicit mentions cover 40.2\% of tasks; a further 11.7\% have
incidental matches only.

That gives three nested strata, reported in Table~\ref{tab:leak}: issues quoting no full key
path (2{,}108 tasks, 668 pull requests on the co-edited key), issues naming no key file as
code (1{,}442 tasks, 549 pull requests), and the strictest set containing no key stem at all
(1{,}243 tasks, 509 pull requests). We treat the middle one as primary. It removes the
leakage channel while leaving the vocabulary signal in place, and it does not select as
heavily for vague issues as the strictest does.

The result is the same under both real controls, which is what makes it credible. Ranks in
this subsection are over the twenty contenders, the scale Table\ref{tab:leak} uses, on which
\m{path\_issue} is fifth on the full corpus. Under the weak full-path control it barely
moves, falling from 0.434 to 0.407 and holding fifth. Under the primary control it falls to 0.361 and \emph{eleventh}; under the strictest,
to 0.336 and thirteenth. \m{hops\_lines} --- pure structure, no issue text --- rises to fifth
in both, at 0.394 and 0.385, ahead of every pure lexical ordering. Because the middle and the
strictest strata agree, the finding is not an artefact of over-correcting: it does not depend
on discarding incidental vocabulary overlap. On the symbol key nothing moves in any stratum:
\m{ppr\_out\_pl}, \m{ppr\_und\_pl} and \m{hops\_lines} are first, second and third in all
three.

We draw two conclusions and decline a third. First, the lexical advantage on the co-edited
key is substantially an effect of issue reports naming the files that change, in code, in
two tasks out of five. Second, that is not purely an artefact: naming does happen in real
reports, so a tool that exploits it exploits something real, and the honest statement is
conditional rather than universal. What we do \emph{not} conclude --- and what an earlier
draft of this paper did conclude --- is that words beat structure on co-edited files in
general. That claim holds on the full corpus and fails under both of the controls we regard
as meaningful.

What survives is fusion. The three cost-scaled fusions hold the top three places on
co-edited files in every stratum, with \m{rrf\_pprpl\_issue\_path} first at 0.431 in the
primary stratum and \m{rrf\_hops\_issue\_path} first at 0.414 in the strictest. The
comparable leakage figures for the seed itself are 10.3\% full path and 35.1\% stem.

The \m{path\_seed} control separates ``issue text'' from ``path matching'' more directly: the
same path-BM25 ranking, queried by the seed's own path instead of the issue, scores 0.371 on
co-edited files against 0.434 for the issue query ($+0.064$, $p_{\text{Holm}} < 10^{-14}$),
and the two are indistinguishable on symbol files ($+0.018$, $p_{\text{Holm}} = 0.30$). The
gain is attributable to the issue text, and only where the target is what developers changed.

\begin{table}
\caption{AUC and rank (in brackets) under the three nested leakage controls: issues quoting
no full key path, issues naming no key file as code, and the strictest set containing no key
stem at all. Ranks are within the twenty contenders on that stratum.}
\label{tab:leak}
\footnotesize
\setlength{\tabcolsep}{3pt}
\begin{tabular}{lcc@{\hskip 0.9em}cc@{\hskip 0.9em}cc}
\toprule
& \multicolumn{2}{c}{no full path} & \multicolumn{2}{c}{no explicit mention}
& \multicolumn{2}{c}{no path or stem} \\
\cmidrule(r){2-3}\cmidrule(lr){4-5}\cmidrule(l){6-7}
Method & co\_edited & symbol & co\_edited & symbol & co\_edited & symbol \\
\midrule
\rows{tab_leak}
\bottomrule
\end{tabular}
\end{table}

\subsection{Sensitivity: does the reversal survive in large repositories?}
\label{sec:size}

The corpus skews small --- the median repository contributes 131 graph nodes --- so the
reversal has to be shown not to be an artefact of repositories one can almost read whole.
Table~\ref{tab:size} repeats the ranking within quartiles of repository size (quartile
boundaries 73, 131 and 246 nodes; the largest repository has 2{,}487). The median of 131
graph nodes and the median of 286{,}385 tokens per repository are consistent: they imply
roughly 2{,}200 tokens, or some 260 lines, per file, which is an ordinary Python module.

The pattern of Table~
ef{tab:main} holds in every quartile, so neither it nor the
qualifications of Sections~
ef{sec:shared} and~
ef{sec:leak} are artefacts of repository
size. It does not widen monotonically with size: the
co-edited-minus-symbol rank gap for \m{ppr\_und\_pl} runs 9, 5, 9, 10 across the quartiles
and for \m{hops\_lines} 5, 1, 3, 3, so the right claim is stability, not a trend. In the
largest quartile, issue-text path matching is fifth on co-edited files and twelfth on symbol
files, while cost-scaled personalised PageRank is eleventh and first respectively. The
query-independent prior degrades fastest of all: on co-edited files in the largest
quartile it scores 0.113 against a random floor of 0.089, which is to say that in the
repositories where retrieval matters most, knowing only which files are central is worth
almost nothing. \m{rrf\_hops\_path} never falls below fifth in any of the eight cells,
and is third or better in five of them.

\begin{table}
\caption{Mean AUC (rank among the twenty contenders) by quartile of repository size in
graph nodes. Q1 $\leq 73$, Q2 $\leq 131$, Q3 $\leq 246$, Q4 up to 2{,}487.}
\label{tab:size}
\footnotesize
\setlength{\tabcolsep}{4pt}
\begin{tabular}{lcccc@{\hskip 1.5em}cccc}
\toprule
& \multicolumn{4}{c}{\m{co\_edited}} & \multicolumn{4}{c}{\m{symbol}} \\
\cmidrule(r){2-5}\cmidrule(l){6-9}
Method & Q1 & Q2 & Q3 & Q4 & Q1 & Q2 & Q3 & Q4 \\
\midrule
\rows{tab_size}
\bottomrule
\end{tabular}
\end{table}

\section{Discussion}
\label{sec:discussion}

\paragraph{What we can and cannot say about the two keys.} The tidy version of this paper
would be that the import graph answers the compiler's question and the issue text answers
the developer's, and that the field has been grading the first against the second. Our own
controls will not support the tidy version, and it is worth being exact about what is left.

On the symbol key, structural methods win on the full corpus and are indistinguishable from
issue-to-path matching once the population is matched to the co-edited key
(Section~\ref{sec:shared}). That key is also built by following imports, so we would have
discounted a win there in any case. On the co-edited key, issue-to-path matching wins on the
full corpus and loses to pure structure once tasks whose issues name a key file in code are
removed (Section~\ref{sec:leak}). Each half of the reversal, in other words, has an
explanation that is not about the signals at all --- one is a sampling artefact, the other a
leakage artefact --- and an earlier draft of this paper reported both halves as findings.

What remains is weaker but real. The two keys still induce different rankings, and a method
tuned on one is not thereby good on the other. But the honest summary of the single-signal
comparison is that neither family dominates once both confounds are controlled, and the
interesting result is elsewhere.

\paragraph{Fusion is the result.} It is the only claim in this paper that survives every
control we applied. Reciprocal-rank fusion of a cost-scaled walk with issue-to-path matching
is the best non-oracle ordering under both keys on the matched population, where it beats
the best pure lexical ordering by 0.055 and 0.067 and the best pure structural one by 0.097
and 0.039 at the repository level, and it holds the top three places on co-edited files in
all three leakage strata. It needs no model,
no index beyond the file list, and no extra hop through the graph.

Two caveats keep this from being stronger than it is. The fusions are the only methods given
both the seed and the issue, so part of their advantage is simply more input rather than
genuine complementarity; the overlap analysis (Section~\ref{sec:rq4}) argues for
complementarity, but a fusion of structure with a non-issue lexical signal would separate the
two cleanly and we have not run one. And the fusion family and the worst-rank criterion were
both added after the design freeze (Section~\ref{sec:design-freeze}), so they are
exploratory, and a held-out split of repositories would be the right way to confirm them.

\paragraph{The field's default key, restated carefully.} Every graph-based localization
system we could check (Table~\ref{tab:priorkeys}) is graded against the files the gold patch
touched --- our \m{co\_edited} source. On that source the best purely structural ordering
ranks tenth of the twenty-three orderings we score, behind matching issue words against file
paths, a method with no graph, no model and no seed. That gap narrows and then reverses once
filename leakage is removed, which is itself the point: the comparison those systems are
implicitly making depends on a property of the issue corpus that none of them report.

Two further qualifications matter. Our task hands every structural method a correct seed file
for free, whereas LocAgent, CoSIL and the rest must localize from the issue text alone; their
task is harder, and a component that underperforms given a free seed might still earn its
place without one. And they are LLM agents that also read the issue, so the graph is one
component among several. We do not claim their advertised component fails to do its work.
The narrower claim we will defend is that on the ground truth those systems themselves
adopt, graph proximity is not obviously the strongest available signal even when handed the
answer's neighbourhood, and that nobody has reported the comparison or the leakage rate.

\paragraph{Filename leakage deserves to be a standard control.} Two tasks in five in this
corpus have an issue that names a key file as code. That is not a benchmark defect --- real
reports really do contain stack traces --- but it means any lexical baseline evaluated
without the stratification is reporting a mixture of two very different regimes, and any
structural method compared against such a baseline is being judged against a moving target.
The control costs a regular expression over text the benchmark already ships. We think
retrieval papers should report the rate and the stratified result as a matter of course, in
the same way they report the corpus size.

\paragraph{What ``within two hops'' is worth.} The two-hop neighbourhood earns its
reputation on recall and loses it on precision: it contains seven of every eight needed
files and also half the repository. Reporting the recall without the size is, on this
corpus, reporting half a statistic.

\paragraph{What this means for an agent's context, and what it does not.} The budget framing
is not decoration: at a median 286{,}385 tokens, most of these repositories do not fit in the
windows current systems run with, so selection is forced rather than optional. Two numbers
bound what selection can achieve. The answer is a median 2.98\% of the repository, so a
well-chosen context could be some 34$\times$ smaller than the whole; and an oracle ordering
fits the complete answer inside 8{,}000 lines for 97.5\% of pull requests where our best
method manages 58.5\%.

We are deliberately not claiming that better selection raises repair rates. We did not run an
agent, and a retrieval improvement need not survive contact with one. What the ceiling
licenses is the weaker statement: a file absent from the context cannot be used, so the share
of tasks whose answer fits the budget bounds any system operating under it. It does not
license the claim that the remaining gap is reachable from the signals we have, only that it
is not closed off by the tasks themselves. Section~\ref{sec:protocol} specifies the
experiment that would convert the bound into a performance claim, in advance of running it.

\paragraph{A note on order versus selection.} An ordering plus a budget is a selection, and
it is the selection that matters here. We make no claim about the position of a file
\emph{within} a prompt, which is a different question with its own literature; every result
in this paper concerns which files are read at all.

\paragraph{The human reading of the same result.} The budget for a developer is attention
rather than tokens, and the results have a plain reading there. Asked to fix the login page,
a developer thinks of the authentication files partly because the words match and partly
because the imports connect; our fusion result says those two routes are largely
complementary, and our leakage result says that how much the words help depends on how
explicitly the report was written. We report the token framing because it makes the budget
concrete and measurable; we do not think the finding is confined to machines.

\paragraph{Recommendations.} (1) Report retrieval results separately for at least one
structural and one behavioural notion of relevance, and report the worst case across them.
(2) Report, and stratify by, whether the query text names the answer files in code.
(3) Match populations before comparing across keys; a key that only exists for multi-file
changes does not cover the same tasks as one that does not.
(4) Treat a query-independent repository prior, not random order, as the floor for
structural methods --- on symbol-style keys it captures 84\% of a personalised walk.
(5) Where a graph-based tool orders context by structure alone, add a lexical signal from
the issue text and cost-scale the structural side before fusing; where one already uses
both, report them separately as well as combined, since on our corpus the combination
beats either part under both keys. (6) Report cost in tokens and report the
ceiling alongside the score, so a reader can tell a hard task from a poor ordering.
(7) Cluster paired tests by repository; pull requests within a repository are not
independent, and ignoring that inflates significance by many orders of magnitude.

\section{A Pre-Registered Protocol for the Downstream Test}
\label{sec:protocol}

This paper measures retrieval, not repair. The obvious next question --- whether a context
assembled by one of these orderings lets a model actually fix the issue --- requires running
models, which we have not done. Rather than gesture at it as future work, we specify it here
in enough detail to be executed or contradicted, in the same spirit as the design freeze
that preceded this study (\S\ref{sec:design-freeze}). We commit to reporting the result
whichever way it falls.

\paragraph{Corpus.} The 1{,}162 scored pull requests of this study. Every one carries a
test specification from the benchmark suite it came from, naming the tests that must change
from failing to passing and those that must keep passing, so resolution can be decided by
execution rather than by patch similarity. We note that only 15 of them overlap SWE-bench Verified, so this is not a
re-run of an existing leaderboard.

\paragraph{Conditions.} Six contexts per task, each truncated to the same token budget:
(i) the seed file alone, a floor; (ii) seed plus the best structural ordering
(\m{ppr\_und\_pl} on the symbol key, \m{hops\_lines} on the co-edited
key --- no single structural ordering is best under both, so the condition is defined per
key); (iii) seed plus the best lexical ordering (\m{path\_issue});
(iv) seed plus the best fusion by worst-case rank (\m{rrf\_hops\_path}); (v) seed plus the true key, a
ceiling; (vi) the repository truncated to the budget in a fixed canonical order, which is
what a system with no retrieval step does. Budgets of 16{,}000 and 64{,}000 tokens, chosen
from Table~\ref{tab:ceiling} to sit either side of the point where most answers begin to
fit.

\paragraph{Models.} Two, of clearly different capability, so that the result is not a
property of one model. A weak model that fails everywhere and a strong model that succeeds
everywhere would both be uninformative; the pilot below exists to detect that before the
full run.

\paragraph{Generation.} A single API call per task and condition, with no tools and no
retrieval available to the model. This is essential rather than incidental: a model that can
search the repository defeats the budget, and the quantity under study is what a fixed
context supports.

\paragraph{Outcome and analysis.} A task is resolved if the generated patch applies and the
specified tests pass. The comparison is paired across conditions on the same tasks, tested
with McNemar's test, Holm-corrected across the condition pairs, with the pull request as the
unit as elsewhere in this paper. We do not require the models to be reliable; we require the
\emph{difference between conditions} to be measurable, and the pairing makes the stochastic
component common to all conditions.

\paragraph{Pilot and stopping rule.} A 150-task pilot on the weaker model precedes the full
run. If the oracle condition resolves under 10\% or over 90\% of pilot tasks, the budget is
outside the informative range and is adjusted before the main run --- and that adjustment is
reported.

\paragraph{What would falsify the implication.} If conditions (ii), (iii) and (iv) do not
differ from (vi) at equal budget, then retrieval quality does not transfer to repair at
these budgets, and the results of this paper should be read as being about retrieval alone.
That outcome is publishable and we would report it.

\section{Threats to Validity}

\paragraph{Construct.} Neither evidence source is ground truth about what a developer needed
to read; they are recoverable proxies with opposite biases, which is why we report both and
never merge them. Coverage per line assumes whole files are read, which overstates the cost
for agents that read selectively. The execution-based source is absent (\S3.3) and would be
the natural third leg.

\paragraph{Internal.} The graph builder resolves imports statically and misses dynamic
imports and plugin registration. The symbol key is produced by a static analyser and
inherits its resolution failures. Both biases favour structural methods. Issue-path leakage is
measured and does substantially drive the co-edited result: it is reported in
Section~\ref{sec:leak} as a finding rather than a reassurance, and the corresponding claim
is stated conditionally throughout. Our leakage classifier is a regular expression over
issue text and will mislabel some mentions in either direction; we report three nested
strata rather than one so that the conclusion does not rest on a single cut.

\paragraph{External.} Python only, file granularity only, and pull requests drawn from
benchmark suites that select for reproducible test-verified fixes --- a population plausibly
better documented than average, which may flatter the issue-text methods. The median
repository here has 131 graph files and very large repositories are under-represented. The
size split in Section~\ref{sec:size} shows the pattern to be stable across size quartiles
rather than widening with size, so we cannot argue that the bias runs in a harmless
direction; we can only say that it does not appear to be a size effect.

\paragraph{Conclusion.} Pull-request-level aggregation, repository-level paired tests, Holm
correction within comparison families and bootstrapping over repositories guard against the
main inflation risks. Effect sizes accompany every test because several differences are
significant and small. The fusion family, the worst-rank criterion and the directed walks
were all added after the design freeze and are exploratory; we have not held out a split of
repositories to confirm them, which is the main remaining weakness of the fusion claim.

\section{Conclusion}

We compared import-graph proximity and issue-text retrieval as reading orders for Python
change tasks on 2{,}412 tasks from 1{,}162 pull requests across 112 repositories, under two
evidence-based definitions of which files a change required. On the full corpus the two
definitions give opposite answers, structural orderings winning on the files a change
references and lexical orderings on the files a developer changed alongside it. Most of
this paper is then about why that headline does not survive contact with two controls.
Matching the populations, since a co-edited key only exists for a multi-file change,
removes the structural advantage on the symbol key. Removing the two tasks in five whose
issue names a key file in code removes the lexical advantage on the co-edited key. Each
control dissolves one half, and neither family dominates once both are applied.

What survives is the combination. Reciprocal-rank fusion of a cost-scaled walk with
issue-text-to-path matching is the best non-oracle ordering under both keys on a matched
population, holds the top places under every leakage stratum, and costs a BM25 index over
the file list. The graph's two-hop neighbourhood, the standard justification for
structural retrieval, turns out to be half the repository, and an oracle fits the complete
answer in 8{,}000 lines for 97.5\% of pull requests where our best method reaches 58.5\%.

We recommend that context-retrieval work for code change report per evidence source,
report and stratify by whether the query names the answer, match populations before
comparing across keys, cluster its tests by repository, use a repository prior rather than
random order as its floor, and stop treating the dependency graph as the only signal in
the room.

\section*{Data Availability}

The replication package contains the import graphs at every base commit, the task
definitions and evidence keys, the issue-token bags, the per-task orderings and coverage
curves, the per-file token counts, and the analysis scripts that produce every number and
figure in this paper. It will be archived at a permanent identifier on acceptance.

\section*{Use of AI Tools}

Claude (Anthropic) was used to write the analysis and method code, to run the experiments
and to draft this manuscript. All numeric results are produced by scripts in the replication
package and were copied mechanically from their output files; the author verified the
design, the claims and the text. Bibliographic entries were drafted with AI assistance and
then checked against arXiv, the ACL Anthology and the publisher listings on 2026-09-23; the
evaluation section of every system in Table~\ref{tab:priorkeys} was read directly rather
than summarised from an abstract. That check corrected eight entries drafted from memory,
including a journal name, an author order, three titles changed between preprint versions,
and one volume recorded as authored that is in fact edited. The corrections are listed in
the replication package.

\bibliographystyle{ACM-Reference-Format}
\bibliography{references}

\end{document}
